\SourceLeaf{251}%
\Needspace{7\baselineskip}
\begin{center}
\textsc{Gauss} an \textsc{Gerling.} Göttingen, 17. April 1844.
\end{center}

Die Art, wie Sie die Volumengleichheit bloss symmetrischer, nicht zugleich
congruenter, Körper beweisen, hat mir viel Vergnügen gemacht. Man könnte
den Nerv davon so hervorheben, dass man sagt,

1) dass man auf diejenigen Pyramiden aufmerksam macht, deren symmetrische
Gegenstücke mit ihnen zugleich congruent sind, welches nemlich diejenigen
sind, an denen zwei Seitenflächen zu einer dritten normal sind und in ihren
Durchschnitten mit dieser gleiche Kanten geben, und

2) dass man nachweist, wie jede Pyramide in 12 Pyramiden von der eben
bezeichneten Art zerlegt werden kann.

Ob die Beweisart neu ist, kann ich übrigens mit Gewissheit nicht verbürgen.
Sehen Sie wenigstens \textsc{Legendres} Geometrie nach, worin versucht ist,
mehreres strenger oder einfacher als sonst geschehen, zu beweisen. Mir
selbst ist in diesem Augenblick das Buch, welches ich nicht selbst besitze,
nicht zur Hand.

Mein \textls{Bedauern} muss ich nun, da jener Satz nicht mehr davon getroffen ist,
auf die andern Sätze der Stereometrie beschränken, die annoch von der
Exhaustionsmethode abhängig sind wie Euklid XII. 5. Vielleicht ist auch
hier noch manches zu verbessern; in diesem Augenblick habe ich nicht Zeit,
dem Gegenstande weiteres Nachdenken zu widmen. \mbox{.\ .\ .\ .\ .\ .}

\Needspace{7\baselineskip}
\begin{center}
[\textsc{Gerling} \emph{an} \textsc{Gauss}. \emph{Marburg, 7. Juli 1844.}]
\end{center}

\{\mbox{.\ .\ .\ .\ .\ .\ .} Was zuerst die Symmetrie betrifft, wo ich im vorletzten Brief
nur die »Thatsache« angab; so scheint mir der \textls{Grund} in unserer Anschauung
der geraden Linie und Ebene zu liegen, vermöge deren wir in der Ebene auf
einer geraden Linie nur \textls{ein} Perpendikel in einem Punkt denken können,
dieses aber nach \textls{beiden Seiten} hin ins Unendliche verlängert, und eben so
